% !TeX root = ../../Skript.tex
\cohead{\Large\textbf{Wurzelgleichungen}}
\fakesubsection{Wurzelgleichungen}
\begin{tcolorbox}[height=3cm]
	\textcolor{loestc}{Unter einer Wurzelgleichung versteht man eine Gleichung in der die gesuchte Variable auch unter einer Wurzel vorkommt.}
\end{tcolorbox}
Will man die Nullstelle einer Wurzelfunktion bestimmen, so muss man eine Wurzelgleichung lösen:

\medskip

\begin{minipage}{\textwidth}
	\adjustbox{valign=t, padding=0ex 0ex 2ex 0ex}{\begin{minipage}{0.5\textwidth-2ex}
					\textbf{Allgemeines Vorgehen}
					
					1. Wurzel isolieren
	\end{minipage}}%
	\adjustbox{valign=t, padding=2ex 0ex 0ex 0ex}{\begin{minipage}{0.4\textwidth-2ex}
					\textbf{Beispiel:} \(\sqrt{x+9}-5=0\)
					\begin{align*}
						\textcolor{loes}{\sqrt{x+9}-5}&\textcolor{loes}{=0\ \vert\ +5}\\
						\textcolor{loes}{\sqrt{x+9}}&\textcolor{loes}{=5}
					\end{align*}
					
	\end{minipage}}%
\end{minipage}

\vspace{1.3cm}

\begin{minipage}{\textwidth}
	\adjustbox{valign=t, padding=0ex 0ex 2ex 0ex}{\begin{minipage}{0.5\textwidth-2ex}
			2. Wurzel durch Potenzieren auflösen
	\end{minipage}}%
	\adjustbox{valign=t, padding=2ex 0ex 0ex 0ex}{\begin{minipage}{0.4\textwidth-2ex}
			\begin{align*}
				\textcolor{loes}{\sqrt{x+9}}&\textcolor{loes}{=5\ \vert\ ^2}\\
				\textcolor{loes}{x+9}&\textcolor{loes}{=25}
			\end{align*}
			
	\end{minipage}}%
\end{minipage}

\vspace{1.3cm}

\begin{minipage}{\textwidth}
\adjustbox{valign=t, padding=0ex 0ex 2ex 0ex}{\begin{minipage}{0.5\textwidth-2ex}
		3. Gleichung lösen
\end{minipage}}%
\adjustbox{valign=t, padding=2ex 0ex 0ex 0ex}{\begin{minipage}{0.4\textwidth-2ex}
		\begin{align*}
			\textcolor{loes}{x+9}&\textcolor{loes}{=25 \vert\ -9}\\
			\textcolor{loes}{x}&\textcolor{loes}{=16}
		\end{align*}
\end{minipage}}%
\end{minipage}

\vspace{1.3cm}

\begin{minipage}{\textwidth}
\adjustbox{valign=t, padding=0ex 0ex 2ex 0ex}{\begin{minipage}{0.5\textwidth-2ex}
		4. Probe
\end{minipage}}%
\adjustbox{valign=t, padding=2ex 0ex 0ex 0ex}{\begin{minipage}{0.4\textwidth-2ex}
		\begin{align*}
			\textcolor{loes}{\sqrt{16+9}-5=\sqrt{25}-5=5-5=0\ \Checkmark}
		\end{align*}
\end{minipage}}%
\end{minipage}

\medskip

Die Prüfung ist wichtig, da das Potenzieren beider Seiten keine Äquivalenzumformung ist, wie folgendes Beispiel zeigt:

\medskip

\begin{minipage}{\textwidth}
	\adjustbox{valign=t, padding=0ex 0ex 2ex 0ex}{\begin{minipage}{0.5\textwidth-2ex}
		
			1. Wurzel isolieren
	\end{minipage}}%
	\adjustbox{valign=t, padding=2ex 0ex 0ex 0ex}{\begin{minipage}{0.4\textwidth-2ex}
			\begin{align*}
				\sqrt{x+9}+5&=0\textcolor{loes}{\ \vert\ -5}\\
				\textcolor{loes}{\sqrt{x+9}}&\textcolor{loes}{=-5}
			\end{align*}
			
	\end{minipage}}%
\end{minipage}

\vspace{1.3cm}

\begin{minipage}{\textwidth}
	\adjustbox{valign=t, padding=0ex 0ex 2ex 0ex}{\begin{minipage}{0.5\textwidth-2ex}
			2. Wurzel durch Potenzieren auflösen
	\end{minipage}}%
	\adjustbox{valign=t, padding=2ex 0ex 0ex 0ex}{\begin{minipage}{0.4\textwidth-2ex}
			\begin{align*}
				\textcolor{loes}{\sqrt{x+9}}&\textcolor{loes}{=-5\ \vert\ ^2}\\
				\textcolor{loes}{x+9}&\textcolor{loes}{=25}
			\end{align*}
			
	\end{minipage}}%
\end{minipage}

\vspace{1.3cm}

\begin{minipage}{\textwidth}
	\adjustbox{valign=t, padding=0ex 0ex 2ex 0ex}{\begin{minipage}{0.5\textwidth-2ex}
			3. Gleichung lösen
	\end{minipage}}%
	\adjustbox{valign=t, padding=2ex 0ex 0ex 0ex}{\begin{minipage}{0.4\textwidth-2ex}
			\begin{align*}
				\textcolor{loes}{x+9}&\textcolor{loes}{=25 \vert\ -9}\\
				\textcolor{loes}{x}&\textcolor{loes}{=16}
			\end{align*}
	\end{minipage}}%
\end{minipage}

\vspace{1.3cm}

\begin{minipage}{\textwidth}
	\adjustbox{valign=t, padding=0ex 0ex 2ex 0ex}{\begin{minipage}{0.5\textwidth-2ex}
			4. Probe
	\end{minipage}}%
	\adjustbox{valign=t, padding=2ex 0ex 0ex 0ex}{\begin{minipage}{0.4\textwidth-2ex}
			\begin{align*}
				\textcolor{loes}{\sqrt{16+9}+5=\sqrt{25}+5=5+5=10=0\ \Lightning}
			\end{align*}
	\end{minipage}}%
\end{minipage}
\textcolor{loes}{Diese Gleichung hat also keine Lösung.}
\newpage
\begin{Exercise}[title={Löse folgende Gleichungen.}, label=wurzelgleichungenA1]
	\begin{enumerate}[label=\alph*)]
		\item \(\sqrt{x+1}=3\)
		\item \(-\sqrt{x-2}=4\)
		\item \(-\sqrt{x-5}=-1\)
		\item \(-\sqrt{x-7}+3=0\)
		\item \(\sqrt{x+4}+1=0\)
		\item \(-\sqrt{x-0,5}-5=-6\)
		\item \(\sqrt{x-\frac{3}{4}}+1=4\)
		\item \(-\sqrt{x+3,7}-1,2=-3,8\)
		\item \(-\sqrt{x+10}-\frac{3}{10}=0\)
		\item \(\sqrt{x-\frac{7}{9}}-3=\frac{2}{3}\)
	\end{enumerate}
\end{Exercise}
\begin{Answer}[ref=wurzelgleichungenA1]
	\begin{enumerate}[label=\alph*)]
		\item \(x=8\)
		\item keine Lösung
		\item \(x=6\)
		\item \(x=16\)
		\item keine Lösung
		\item \(x=1.5\)
		\item \(x=9,75\)
		\item \(x=3,06\)
		\item keine Lösung
		\item \(x=\frac{128}{9}\)
	\end{enumerate}
\end{Answer}

